\chapter{Research Methods for Interactive Systems Design}

When designing interactive systems, it is crucial to understand the design space:
\begin{itemize}
    \item Users: needs and requirements
    \item Situation: what is the current state of the problem
    \item Context: location, legislation, culture, etc.
\end{itemize}

There are different types of research methods:
\begin{itemize}
    \item \ul{Quantitative / Qualitative}: Quantitative methods involve collecting and analyzing numerical data, while qualitative methods focus on understanding human behavior and the reasons behind it.
    \item \ul{In situ / In Laboratory}: In situ research is conducted in the natural environment of the users, while laboratory research takes place in a controlled setting.
    \item \ul{Unstructured / Structured}: Unstructured research allows for more flexibility and exploration, while structured research follows a predefined plan and methodology.
\end{itemize}

In order to design effective interactive systems, \emph{triangulation} is often used, which involves combining multiple research methods to gain a comprehensive understanding of the design space. During the following sections, we will explore various research methods that can be employed to gather all the needed information.

\section{Ethnography}

Ethnography is a research method that is founded on participation rather than observation, with the goal of not only understanding what people do but also why they do it and how they feel about it. It involves immersing oneself in the users' environment to gain insights into their behaviors, motivations, and experiences.
\begin{itemize}
    \item It is a qualitative, in situ and unstructured research method that uses induction.
\end{itemize}

There are two main types of ethnography:
\begin{description}
    \item[Observing Participation] The researcher observes the users' activities and interactions without actively engaging in them. This allows for an objective perspective on the users' behaviors and experiences.
    \item[Participating Observation] The researcher is actively involved in the users' activities and interactions, allowing for a deeper understanding of their experiences and perspectives. This method has the risk of going ``native'', where the researcher becomes too immersed in the users' environment and loses objectivity.
\end{description}

Some methods used in ethnography include:
\begin{itemize}
    \item \ul{Field Diary}: Personal notes taken by the researcher during their immersion in the field, capturing observations, thoughts, and reflections.
    \item \ul{Cameras}: Used to document the environment and interactions, providing visual data for analysis. One important disadvantage is the \emph{Hawthorne Effect}, which refers to the phenomenon where individuals modify their behavior in response to being observed, potentially leading to biased data.
    \item \ul{Diary Studies}: Participants keep a diary of their activities and experiences.
    \item \ul{People receive package with material for documentation}: Participants are given materials to document their experiences, such as cameras or notebooks. After a certain period, the researcher collects the materials and analyzes the data to gain insights into the users' behaviors and experiences.
\end{itemize}



\section{Action Research}

AR is a research method that is based on the idea of taking action to solve a problem while simultaneously conducting research to understand the problem and its context. As opposed to ethnography (where the researcher is an insider but does not change the environment), in AR the researcher is an insider and changes the environment.
\begin{itemize}
    \item It is a qualitative, in situ and unstructured research method. It uses induction, but its generalization is difficult because it is context-specific.
\end{itemize}

There are three main types of AR:
\begin{description}
    \item[Technical AR] The researcher intervenes based on his technical expertise, sharing his knowledge with the users to solve a specific problem.
    \item[Practical-deliverative AR] The researcher tries to understand the users' needs and requirements, and then tries to solve the problem based on that understanding (not necessarily based on his previous technical expertise).
    \item[Critical-emancipatory AR] The researcher tries to understand the users' needs and requirements, and then tries to empower the users to solve the problem by themselves, without necessarily providing a solution directly.
\end{description}

AR involves a cyclical process of planning, acting and examining the results, with feedback loops that allow for continuous improvement and adaptation of the interventions. The three main steps of AR are:
\begin{enumerate}
    \item \ul{Planning}: The researcher identifies a problem and develops a plan to address it. \emph{``Unfreezing''} the current state of the problem is crucial to allow for change, as it involves breaking down existing habits and behaviors that may be resistant to change.
    \item \ul{Acting}: The researcher implements the plan and takes action to address the problem. This may involve introducing new technologies, processes, or interventions to improve the situation.
    \item \ul{Examining the results}: The researcher evaluates the outcomes of the actions taken and reflects on the effectiveness of the interventions. \emph{``Refreezing''} the new state of the problem is important to ensure that the changes are sustained over time, as it involves reinforcing new habits and behaviors that have been established.
\end{enumerate}

\section{Research through Design}

RtD is a research method that is based on the idea of using design as a means of inquiry and exploration. It involves creating prototypes and artifacts not to solve a specific problem, but to gain insights and understanding about the design space and the users' needs and requirements.
\begin{itemize}
    \item It is a qualitative, in situ and usually unstructured research method.
\end{itemize}

It is a cyclical process that involves four main steps:
\begin{enumerate}
    \item \ul{Selecting a problem}: The researcher identifies a problem or area of interest to explore through design.
    \item \ul{Designing a solution}: The researcher creates a prototype or artifact that addresses the selected problem or area of interest.
    \item \ul{Evaluating the solution}: The researcher evaluates the prototype or artifact to gain insights into the design space and the users' needs and requirements.
    \item \ul{Reflecting on the results}: The researcher reflects on the insights gained from the evaluation and uses them to inform the next iteration of the design process.
\end{enumerate}

It has several benefits, such as:
\begin{itemize}
    \item It shows technical opportunities and challenges that engineers may not have thought of.
    \item Exposes the gaps between the theory and practice of design.
    \item Changes current practices, therefore opening up new design spaces.
    \item Reveals designer patterns and assumptions, which can be used to inform future design decisions.
\end{itemize}

\section{Field Deployment}

Field deployment is a research method that involves deploying a prototype or artifact in the users' natural environment to gain insights into its use and effectiveness. It allows researchers to observe how users interact with the prototype in real-world settings, providing valuable feedback for further design iterations. Testing it in the lab is not enough, as it may not capture the complexities and nuances of real-world use.
\begin{itemize}
    \item It is a qualitative, in situ and usually structured research method.
\end{itemize}



There are 7 main steps in field deployment:
\begin{enumerate}
    \item \ul{Finding a Setting}: There are three different types of settings:
    \begin{itemize}
        \item \emph{Convenience}: Deployment within the researcher's own environment. Not representative, and participants may be biased.
        \item \emph{Semi-Controlled}: It involves knwon and unknown participants. It can be more representative, but it may still have some biases.
        \item \emph{In the Wild}: Deployment in the users' natural environment to almost unknown participants. It is the most representative, but needs more robust prototypes.
    \end{itemize}
    \item \ul{Defining Goals}: The fundamental goal is increasing understanding of the prototype that generalizes beyond the specific deployment.
    \item \ul{Recruiting Participants}: They should be representative of the target user population, also taking into account the researcher's group.
    \item \ul{Designing Data Collection Instruments}: Both qualitative (e.g., interviews, observations) and quantitative (e.g., measured time to complete a task, error rates) data collection instruments should be designed to gather comprehensive insights into the users' interactions with the prototype.
    \item \ul{Conducting the Deployment}: Pilot tests should be conducted to identify and address any issues with the prototype or data collection instruments before the full deployment.
    \item \ul{Ending the Deployment}: When the deployment is over, the product is usually taken away from the users in order to reduce the cost of mantaining its servers. In addition, they are usually not allowed to be bought by the users, as usually researchers do not want to sell products but only to gain insights.
    
    Taking the prototype away from the users can be a difficult process, as it may have become an integral part of their daily lives. It is important to plan for this transition and provide support to users as they adjust to the change.
    \item \ul{Analyzing the Data}
\end{enumerate}


\section{Experimental Research}

Experimental research is a research method whose aim is to know the cause-effect relationship between two or more variables. It involves manipulating one variable (called \emph{independent variable}) and observing the effect on one or more other variables (called \emph{dependent variables}), while controlling for other factors that may influence the results. In order to conduct a good experiment, a prediction (called \emph{hypothesis}) should be made about the expected relationship between the variables.
\begin{itemize}
    \item It is a quantitative, in laboratory and structured research method.
\end{itemize}

Two different factors should be taken into account when designing an experiment:
\begin{itemize}
    \item \ul{Validity}: It refers to the extent to which the results of an experiment can be generalized to other settings, populations, and times. There are several types of validity, as population validity, ecological validity, temporal validity, etc.
    \item \ul{Reliability}: It refers to the consistency and stability of the results of an experiment. An experiment is considered reliable if it produces the same results when repeated under the same conditions.
\end{itemize}

As explained before, hypotheses are predictions about the expected relationship between variables, and they are essential for designing and conducting experiments. They should be \emph{precise}, \emph{meaningful} (based on theory), \emph{testable}, and \emph{falsifiable}. In order to test a hypothesis, the null hypothesis (which states that there is no relationship between the variables) should be rejected.

There are two different approaches to research in experimental research:
\begin{itemize}
    \item \emph{Between-Subject Design}: Different groups of participants are exposed to different conditions of the independent variable. This design allows for comparisons between groups, but it may require a larger sample size to achieve statistical power.
    \item \emph{Within-Subject Design}: The same group of participants is exposed to all conditions of the independent variable. They are firstly exposed to one condition, and then to the other condition. This design allows for comparisons within the same group, but it may be affected by order effects (e.g., fatigue, learning).
    
    In order to mitigate order effects in within-subject designs, counterbalancing can be used, which involves varying the order of conditions across participants to ensure that any potential order effects are evenly distributed and do not bias the results. There are two main types of counterbalancing:
    \begin{itemize}
        \item \emph{Complete Counterbalancing}: All possible orders of conditions are used, which can be impractical for a large number of conditions, as $x!$ orders are needed (where $x$ is the number of conditions).
        \item \emph{Latin Square Counterbalancing}: A more efficient method that uses a Latin square design to ensure that each condition appears in each position an equal number of times, reducing the number of orders needed to $x$ (where $x$ is the number of conditions).
    \end{itemize}
\end{itemize}


% // TODO: Ver vídeo. Explica Measures applied?


\section{Survey Research}

Survey research is a research method that involves collecting data from a sample of individuals using standardized questionnaires or interviews.
\begin{itemize}
    \item It is a quantitative, both in situ and in laboratory, and a structured research method.
\end{itemize}

Surveys are used to gather information about people's attitudes, beliefs, behaviors, and characteristics... but are not good when understanding the reasons behind those attitudes, beliefs, behaviors, and characteristics. Therefore, they are often used in combination with other research methods (e.g., ethnography, action research) to gain a more comprehensive understanding of the design space.
\begin{itemize}
    \item Surveys are used to decide if the results of the research methods can be generalized to a larger population.
    \item The other research methods are used to understand the reasons or trends behind the attitudes, beliefs, behaviors, and characteristics of the users.
\end{itemize}

There are different stages in survey research:
\begin{enumerate}
    \item \ul{Research Goals and Constructs}:
    Firstly, the research goals should be defined, which will guide the entire survey research process. Then, the constructs that will be measured should be identified.
    \item \ul{Population and Sampling}: There are different scopes of people that can be surveyed. The population groups considered, from the most general to the most specific, are:
    \begin{itemize}
        \item \emph{Population}: The entire group of individuals that the researcher is interested in studying.
        \item \emph{Sampling Frame}: The individuals that the researcher has access to and can potentially include in the survey.
        \item \emph{Sample}: The individuals that are actually asked to participate in the survey. Not all individuals in the sampling frame may be included in the sample, as some may be excluded due to various reasons (e.g., ethical reasons, ineligibility). There are different sampling methods that can be used to select the sample:
        \begin{itemize}
            \item \emph{Probability Sampling}: Each individual in the sampling frame has the same probability of being selected for the sample. This method allows for generalization of the results to the population, but it may be more time-consuming and expensive.
            \item \emph{Non-Probability Sampling}: Individuals are selected based on non-random criteria, such as convenience or judgment. This method is less time-consuming and expensive, but it may introduce bias and limit the generalizability of the results.
        \end{itemize}
        \item \emph{Respondents}: The individuals that actually respond to the survey. The response rate (i.e., the percentage of respondents out of the total sample) is an important factor to consider, as a low response rate can lead to non-response bias and affect the generalizability of the results.
    \end{itemize}
    \item \ul{Questionnaire Design and Bias}:
    
    There are two main types of questions that can be used in a survey:
    \begin{itemize}
        \item \emph{Close Ended Questions}: These questions provide respondents with a set of predefined response options to choose from. They are easier to analyze and easier for respondents to answer, but they may limit the depth and richness of the data collected. This should be the most used type of question in surveys. They can be further categorized into:
        \begin{itemize}
            \item Single-choice questions.
            \item Multiple-choice questions.
            \item Likert scale questions, aka rating questions.
            \item Ranking questions.
        \end{itemize}
        \item \emph{Open Ended Questions}: These questions allow respondents to provide their own answers in their own words. They can provide richer and more detailed data, but they may be more difficult to analyze and may require more time for respondents to answer.
    \end{itemize}


    When designing a questionnaire, it is important to be aware of potential biases that can affect the validity of the results. A bias is a systematic error that can lead to inaccurate or misleading results. Some common types of bias in survey research include:
    \begin{itemize}
        \item \emph{Satisficing Bias}: Respondents may use a low cognitive effort strategy to answer survey questions, such as selecting the first acceptable response option.
        \item \emph{Acquiescence Bias}: Respondents may have a tendency to agree with statements or questions, regardless of their true feelings or opinions. Questions with \verb|agree/disagree| response options are particularly susceptible to this bias.
        \item \emph{Social Desirability Bias}: Respondents may provide answers that they believe are socially acceptable or desirable, rather than their true thoughts or feelings.
        \item \emph{Response Order Bias}: The order in which response options are presented can influence respondents' choices, with a tendency to select options that are presented first or last.
        \item \emph{Question Order Bias}: The order of questions can influence respondents' answers, as earlier questions may prime certain thoughts or feelings that affect responses to later questions.
    \end{itemize}

    In addition, when designing a questionnaire, there are already existing validated questionnaires that can be used to measure certain constructs. Using these existing questionnaires can save time and effort, and can also improve the validity and reliability of the results, as these questionnaires have already been tested and validated in previous research.
    \item \ul{Review and Survey Pretesting}:
    
    There are two main types of pretesting that can be conducted to evaluate the questionnaire before it is administered to the full sample:
    \begin{itemize}
        \item \emph{Cognitive Pretesting}: This involves asking a small group of respondents to complete the questionnaire while thinking aloud, allowing the researcher to identify any issues with question wording, comprehension, or response options.
        \item \emph{Pilot Testing}: This involves administering the questionnaire to a small sample of respondents that are similar to the target population, allowing the researcher to identify any issues with the survey administration process, such as the length of time it takes to complete the survey or any technical issues with the survey platform.
    \end{itemize}
    \item \ul{Implementation and Launch}:
    
    There are different well known platforms that can be used to implement and launch a survey, such as Qualtrics, SurveyMonkey, Google Forms, etc.
    \item \ul{Data Analysis and Reporting}:
    
    When analyzing survey data, it is important to consider the type of data collected (e.g., nominal, ordinal, interval, ratio) and to use appropriate statistical methods for analysis. Some additional metrics can be analyzed, such as:
    \begin{itemize}
        \item \emph{Data Logging}: Some survey platforms allow for data logging, which involves recording additional information about respondents' interactions with the survey, such as the time taken to complete the survey or the order of responses. This can provide valuable insights into respondents' behavior and can help identify potential issues with the survey design.
        \item \emph{Response Rates}: Analyzing response rates can help identify potential biases in the sample and can inform strategies for improving response rates in future surveys.
    \end{itemize}
\end{enumerate}